% Options for packages loaded elsewhere
\PassOptionsToPackage{unicode}{hyperref}
\PassOptionsToPackage{hyphens}{url}
%
\documentclass[
  12pt,
  a4paper,
  DIV=11,
  numbers=noendperiod]{scrreprt}

\usepackage{amsmath,amssymb}
\usepackage{setspace}
\usepackage{iftex}
\ifPDFTeX
  \usepackage[T1]{fontenc}
  \usepackage[utf8]{inputenc}
  \usepackage{textcomp} % provide euro and other symbols
\else % if luatex or xetex
  \usepackage{unicode-math}
  \defaultfontfeatures{Scale=MatchLowercase}
  \defaultfontfeatures[\rmfamily]{Ligatures=TeX,Scale=1}
\fi
\usepackage{lmodern}
\ifPDFTeX\else  
    % xetex/luatex font selection
\fi
% Use upquote if available, for straight quotes in verbatim environments
\IfFileExists{upquote.sty}{\usepackage{upquote}}{}
\IfFileExists{microtype.sty}{% use microtype if available
  \usepackage[]{microtype}
  \UseMicrotypeSet[protrusion]{basicmath} % disable protrusion for tt fonts
}{}
\usepackage{xcolor}
\setlength{\emergencystretch}{3em} % prevent overfull lines
\setcounter{secnumdepth}{5}
% Make \paragraph and \subparagraph free-standing
\makeatletter
\ifx\paragraph\undefined\else
  \let\oldparagraph\paragraph
  \renewcommand{\paragraph}{
    \@ifstar
      \xxxParagraphStar
      \xxxParagraphNoStar
  }
  \newcommand{\xxxParagraphStar}[1]{\oldparagraph*{#1}\mbox{}}
  \newcommand{\xxxParagraphNoStar}[1]{\oldparagraph{#1}\mbox{}}
\fi
\ifx\subparagraph\undefined\else
  \let\oldsubparagraph\subparagraph
  \renewcommand{\subparagraph}{
    \@ifstar
      \xxxSubParagraphStar
      \xxxSubParagraphNoStar
  }
  \newcommand{\xxxSubParagraphStar}[1]{\oldsubparagraph*{#1}\mbox{}}
  \newcommand{\xxxSubParagraphNoStar}[1]{\oldsubparagraph{#1}\mbox{}}
\fi
\makeatother


\providecommand{\tightlist}{%
  \setlength{\itemsep}{0pt}\setlength{\parskip}{0pt}}\usepackage{longtable,booktabs,array}
\usepackage{calc} % for calculating minipage widths
% Correct order of tables after \paragraph or \subparagraph
\usepackage{etoolbox}
\makeatletter
\patchcmd\longtable{\par}{\if@noskipsec\mbox{}\fi\par}{}{}
\makeatother
% Allow footnotes in longtable head/foot
\IfFileExists{footnotehyper.sty}{\usepackage{footnotehyper}}{\usepackage{footnote}}
\makesavenoteenv{longtable}
\usepackage{graphicx}
\makeatletter
\newsavebox\pandoc@box
\newcommand*\pandocbounded[1]{% scales image to fit in text height/width
  \sbox\pandoc@box{#1}%
  \Gscale@div\@tempa{\textheight}{\dimexpr\ht\pandoc@box+\dp\pandoc@box\relax}%
  \Gscale@div\@tempb{\linewidth}{\wd\pandoc@box}%
  \ifdim\@tempb\p@<\@tempa\p@\let\@tempa\@tempb\fi% select the smaller of both
  \ifdim\@tempa\p@<\p@\scalebox{\@tempa}{\usebox\pandoc@box}%
  \else\usebox{\pandoc@box}%
  \fi%
}
% Set default figure placement to htbp
\def\fps@figure{htbp}
\makeatother

\usepackage[utf8]{inputenc}
\usepackage[T2A]{fontenc}
\usepackage[english,russian]{babel}
\usepackage{fontspec}
\setmainfont{Liberation Serif}
\setsansfont{Liberation Sans}
\setmonofont{Liberation Mono}
\KOMAoption{captions}{tableheading}
\makeatletter
\@ifpackageloaded{caption}{}{\usepackage{caption}}
\AtBeginDocument{%
\ifdefined\contentsname
  \renewcommand*\contentsname{Содержание}
\else
  \newcommand\contentsname{Содержание}
\fi
\ifdefined\listfigurename
  \renewcommand*\listfigurename{Список иллюстраций}
\else
  \newcommand\listfigurename{Список иллюстраций}
\fi
\ifdefined\listtablename
  \renewcommand*\listtablename{Список таблиц}
\else
  \newcommand\listtablename{Список таблиц}
\fi
\ifdefined\figurename
  \renewcommand*\figurename{Рисунок}
\else
  \newcommand\figurename{Рисунок}
\fi
\ifdefined\tablename
  \renewcommand*\tablename{Таблица}
\else
  \newcommand\tablename{Таблица}
\fi
}
\@ifpackageloaded{float}{}{\usepackage{float}}
\floatstyle{ruled}
\@ifundefined{c@chapter}{\newfloat{codelisting}{h}{lop}}{\newfloat{codelisting}{h}{lop}[chapter]}
\floatname{codelisting}{Список}
\newcommand*\listoflistings{\listof{codelisting}{Листинги}}
\makeatother
\makeatletter
\makeatother
\makeatletter
\@ifpackageloaded{caption}{}{\usepackage{caption}}
\@ifpackageloaded{subcaption}{}{\usepackage{subcaption}}
\makeatother

\ifLuaTeX
\usepackage[bidi=basic]{babel}
\else
\usepackage[bidi=default]{babel}
\fi
\babelprovide[main,import]{russian}
\babelprovide[import]{english}
% get rid of language-specific shorthands (see #6817):
\let\LanguageShortHands\languageshorthands
\def\languageshorthands#1{}
\usepackage[style=gost-numeric,backend=biber,langhook=extras,autolang=other*]{biblatex}
\usepackage{csquotes}
\usepackage{bookmark}

\IfFileExists{xurl.sty}{\usepackage{xurl}}{} % add URL line breaks if available
\urlstyle{same} % disable monospaced font for URLs
\hypersetup{
  pdftitle={Лабораторная работа №2},
  pdfauthor={София Черная},
  pdflang={ru-RU},
  hidelinks,
  pdfcreator={LaTeX via pandoc}}


\title{Лабораторная работа №2}
\usepackage{etoolbox}
\makeatletter
\providecommand{\subtitle}[1]{% add subtitle to \maketitle
  \apptocmd{\@title}{\par {\large #1 \par}}{}{}
}
\makeatother
\subtitle{Модель SIR и модель Лотки-Вольтерры}
\author{София Черная}
\date{}

\begin{document}
\maketitle

\renewcommand*\contentsname{Содержание}
{
\setcounter{tocdepth}{1}
\tableofcontents
}
\listoffigures
\listoftables

\setstretch{1.5}
\chapter{Цель
работы}\label{ux446ux435ux43bux44c-ux440ux430ux431ux43eux442ux44b}

Изучение модели SIR (эпидемиологической модели) и модели Лотки-Вольтерры
(системы «хищник-жертва»). Освоение методов имитационного моделирования
систем дифференциальных уравнений с использованием языка Julia и пакета
DrWatson.

\chapter{Задание}\label{ux437ux430ux434ux430ux43dux438ux435}

\begin{enumerate}
\def\labelenumi{\arabic{enumi}.}
\tightlist
\item
  Создать рабочее пространство для второй лабораторной работы.
\item
  Реализовать модель SIR в базовом варианте.
\item
  Преобразовать код в литературный стиль с помощью Literate.jl.
\item
  Сгенерировать из литературного кода чистый код, Jupyter Notebook и
  документацию Quarto.
\item
  Провести параметрическое исследование модели SIR.
\item
  Реализовать модель Лотки-Вольтерры в базовом и параметрическом
  вариантах.
\item
  Провести анализ результатов, оформить графики и включить их в отчёт.
\end{enumerate}

\chapter{Теоретическое
введение}\label{ux442ux435ux43eux440ux435ux442ux438ux447ux435ux441ux43aux43eux435-ux432ux432ux435ux434ux435ux43dux438ux435}

\section{Модель SIR}\label{ux43cux43eux434ux435ux43bux44c-sir}

Модель SIR является классической математической моделью эпидемиологии,
описывающей распространение инфекционного заболевания в замкнутой
популяции.

\subsection{Компартменты
модели}\label{ux43aux43eux43cux43fux430ux440ux442ux43cux435ux43dux442ux44b-ux43cux43eux434ux435ux43bux438}

Модель делит популяцию на три группы: - \textbf{S (Susceptible,
Восприимчивые)} --- индивиды, которые могут заразиться. - \textbf{I
(Infectious, Инфицированные)} --- индивиды, которые больны и могут
передавать инфекцию. - \textbf{R (Recovered, Выздоровевшие)} ---
индивиды, которые переболели и приобрели иммунитет (или умерли). Они
больше не участвуют в распространении болезни.

\subsection{Система дифференциальных
уравнений}\label{ux441ux438ux441ux442ux435ux43cux430-ux434ux438ux444ux444ux435ux440ux435ux43dux446ux438ux430ux43bux44cux43dux44bux445-ux443ux440ux430ux432ux43dux435ux43dux438ux439}

Динамика модели описывается следующей системой уравнений:

\[
\begin{cases}
\dfrac{dS}{dt} = -\beta \dfrac{S I}{N} \\[1em]
\dfrac{dI}{dt} = \beta \dfrac{S I}{N} - \gamma I \\[1em]
\dfrac{dR}{dt} = \gamma I
\end{cases}
\]

где: - ( N = S + I + R ) --- общая (постоянная) численность популяции; -
( \beta ) --- коэффициент заражения (скорость передачи инфекции); - (
\gamma ) --- коэффициент выздоровления (величина, обратная средней
продолжительности болезни ( 1/\gamma )).

\subsection{Ключевые
показатели}\label{ux43aux43bux44eux447ux435ux432ux44bux435-ux43fux43eux43aux430ux437ux430ux442ux435ux43bux438}

\begin{itemize}
\tightlist
\item
  \textbf{Базовое репродуктивное число}: ( R\_0 = \dfrac{\beta}{\gamma}
  ). Показывает, сколько в среднем человек заразит один инфицированный в
  полностью восприимчивой популяции. При ( R\_0 \textgreater{} 1 )
  эпидемия разрастается, при ( R\_0 \textless{} 1 ) --- затухает.
\item
  \textbf{Порог коллективного иммунитета}: ( 1 - \dfrac{1}{R_0} ). Доля
  популяции, которая должна иметь иммунитет для предотвращения роста
  эпидемии.
\end{itemize}

\section{Модель
Лотки-Вольтерры}\label{ux43cux43eux434ux435ux43bux44c-ux43bux43eux442ux43aux438-ux432ux43eux43bux44cux442ux435ux440ux440ux44b}

Модель Лотки-Вольтерры (или модель «хищник-жертва») описывает динамику
взаимодействия двух биологических видов.

\subsection{Система
уравнений}\label{ux441ux438ux441ux442ux435ux43cux430-ux443ux440ux430ux432ux43dux435ux43dux438ux439}

\[
\begin{cases}
\dfrac{dx}{dt} = \alpha x - \beta x y \\[1em]
\dfrac{dy}{dt} = \delta x y - \gamma y
\end{cases}
\]

где: - ( x ) --- численность популяции жертв (например, зайцев); - ( y )
--- численность популяции хищников (например, лис); - ( \alpha ) ---
коэффициент естественного прироста жертв; - ( \beta ) --- коэффициент
смертности жертв в результате встреч с хищниками; - ( \delta ) ---
коэффициент прироста хищников за счёт потребления жертв; - ( \gamma )
--- коэффициент естественной смертности хищников.

\subsection{Характеристики
модели}\label{ux445ux430ux440ux430ux43aux442ux435ux440ux438ux441ux442ux438ux43aux438-ux43cux43eux434ux435ux43bux438}

Система имеет две стационарные точки. Нетривиальная точка
сосуществования видов: {[} x\^{}* = \frac{\gamma}{\delta},\quad y\^{}* =
\frac{\alpha}{\beta} {]} Вокруг этой точки решения носят колебательный
характер, причём пик численности хищников отстаёт от пика численности
жертв.

\chapter{Выполнение лабораторной
работы}\label{ux432ux44bux43fux43eux43bux43dux435ux43dux438ux435-ux43bux430ux431ux43eux440ux430ux442ux43eux440ux43dux43eux439-ux440ux430ux431ux43eux442ux44b}

Для выполнения лабораторной работы перехожу в уже созданный репозиторий
labs/lab02.

\section{Модель SIR}\label{ux43cux43eux434ux435ux43bux44c-sir-1}

Создаю файл sir\_ode.jl (рис. Рисунок~\ref{fig-001}).

\begin{figure}

\centering{

\includegraphics[width=0.7\linewidth,height=\textheight,keepaspectratio]{image/1.png}

}

\caption{\label{fig-001}Создание файла через команду touch}

\end{figure}%

Вставляю предложенный из методических материалов программный код в файл
(рис. Рисунок~\ref{fig-002}). Код реализует функцию \texttt{sir\_ode!},
которая описывает правые части системы дифференциальных уравнений, а
также задаёт параметры модели и начальные условия.

\begin{figure}

\centering{

\includegraphics[width=0.7\linewidth,height=\textheight,keepaspectratio]{image/2.png}

}

\caption{\label{fig-002}Программный код модели SIR}

\end{figure}%

Запускаю скрипт с помощью команды
\texttt{julia\ -\/-project=.\ scripts/sir\_ode.jl} (рис.
Рисунок~\ref{fig-003}).

\begin{figure}

\centering{

\includegraphics[width=0.7\linewidth,height=\textheight,keepaspectratio]{image/3.png}

}

\caption{\label{fig-003}Запуск скрипта sir\_ode.jl}

\end{figure}%

Вывод в консоли параметров модели (рис. Рисунок~\ref{fig-004}). Здесь
отображаются заданные значения коэффициентов ( \beta ) и ( \gamma ),
рассчитанное базовое репродуктивное число ( R\_0 ), а также начальные
условия.

\begin{figure}

\centering{

\includegraphics[width=0.7\linewidth,height=\textheight,keepaspectratio]{image/4.png}

}

\caption{\label{fig-004}Параметры модели SIR, выведенные в консоль}

\end{figure}%

В результате выполнения скрипта были сгенерированы несколько графиков,
позволяющих проанализировать динамику эпидемии.

На рис. Рисунок~\ref{fig-005} представлен график «Динамика эффективного
репродуктивного числа» ( R\_e(t) = R\_0 \cdot S(t)/N ). - При ( R\_e(t)
\textgreater{} 1 ): количество новых зараженных растет, эпидемия
развивается. - При ( R\_e(t) = 1 ): система находится в равновесии. Это
момент пика эпидемии, когда число больных ( I(t) ) перестает расти и
начинает падать. - При ( R\_e(t) \textless{} 1 ): количество новых
зараженных падает, эпидемия затухает.

\begin{figure}

\centering{

\includegraphics[width=0.7\linewidth,height=\textheight,keepaspectratio]{image/5.png}

}

\caption{\label{fig-005}Динамика эффективного репродуктивного числа (
R\_e(t) )}

\end{figure}%

Следующим идет график «Динамика числа зараженных» (рис.
Рисунок~\ref{fig-006}). На нем показано изменение числа инфицированных (
I(t) ) во времени. Хорошо виден колоколообразный пик эпидемии.

\begin{figure}

\centering{

\includegraphics[width=0.7\linewidth,height=\textheight,keepaspectratio]{image/6.png}

}

\caption{\label{fig-006}Динамика числа зараженных ( I(t) )}

\end{figure}%

График «Экспоненциальный рост (лог. шкала)» (рис. Рисунок~\ref{fig-007})
показывает ту же кривую ( I(t) ), но в полулогарифмическом масштабе.
Прямолинейный участок в начале графика подтверждает экспоненциальный
характер роста числа зараженных на начальном этапе эпидемии, что
соответствует теоретическим предсказаниям для модели SIR при ( R\_0
\textgreater{} 1 ).

\begin{figure}

\centering{

\includegraphics[width=0.7\linewidth,height=\textheight,keepaspectratio]{image/7.png}

}

\caption{\label{fig-007}Экспоненциальный рост числа зараженных в
логарифмическом масштабе}

\end{figure}%

На графике «Модель SIR: динамика эпидемии» (рис. Рисунок~\ref{fig-008})
представлена полная картина: показаны кривые для всех трех групп ---
восприимчивых ( S(t) ), инфицированных ( I(t) ) и выздоровевших ( R(t)
). Видно, как по мере развития эпидемии число восприимчивых падает, а
число выздоровевших растет.

\begin{figure}

\centering{

\includegraphics[width=0.7\linewidth,height=\textheight,keepaspectratio]{image/8.png}

}

\caption{\label{fig-008}Модель SIR: динамика всех трёх групп населения}

\end{figure}%

На рис. Рисунок~\ref{fig-009} показана компактная панель из шести
графиков, объединяющая различные аспекты анализа модели SIR: динамику
популяций, фазовый портрет, скорости изменения и спектральный анализ.

\begin{figure}

\centering{

\includegraphics[width=0.7\linewidth,height=\textheight,keepaspectratio]{image/9.png}

}

\caption{\label{fig-009}Компактная панель графиков для анализа модели
SIR}

\end{figure}%

График «Динамика эпидемии в процентах» (рис. Рисунок~\ref{fig-010})
аналогичен графику на рис. Рисунок~\ref{fig-008}, но значения по оси Y
представлены в процентах от общей численности популяции. Фиолетовая
пунктирная линия обозначает «порог коллективного иммунитета» ( 1-1/R\_0
). Когда доля восприимчивых ( S/N ) опускается ниже этого порога,
эпидемия начинает затухать, так как каждый больной в среднем заражает
меньше одного человека.

\begin{figure}

\centering{

\includegraphics[width=0.7\linewidth,height=\textheight,keepaspectratio]{image/10.png}

}

\caption{\label{fig-010}Динамика эпидемии в процентах от общей
численности популяции}

\end{figure}%

На фазовом портрете (рис. Рисунок~\ref{fig-011}) показана внутренняя
структура процесса. По оси Ox --- число восприимчивых ( S ), по оси Oy
--- число инфицированных ( I ). Траектория системы начинается в точке с
максимальным ( S ) и движется к оси ( S ), где ( I = 0 ), что
соответствует окончанию эпидемии.

\begin{figure}

\centering{

\includegraphics[width=0.7\linewidth,height=\textheight,keepaspectratio]{image/11.png}

}

\caption{\label{fig-011}Фазовый портрет модели SIR}

\end{figure}%

После успешного выполнения базового скрипта, я преобразовал файл
\texttt{sir\_ode.jl} в литературный код, добавив подробные комментарии и
теоретические пояснения. Полученный файл назван
\texttt{sir\_literode.jl} (рис. Рисунок~\ref{fig-012}). Такой подход
соответствует методологии литературного программирования.

\begin{figure}

\centering{

\includegraphics[width=0.7\linewidth,height=\textheight,keepaspectratio]{image/12.png}

}

\caption{\label{fig-012}Литературный код модели SIR в файле
sir\_literode.jl}

\end{figure}%

Запускаю скрипт \texttt{sir\_literode.jl} для генерации производных
форматов (рис. Рисунок~\ref{fig-013}). Для этого используется
специальный скрипт \texttt{tangle.jl}.

\begin{figure}

\centering{

\includegraphics[width=0.7\linewidth,height=\textheight,keepaspectratio]{image/13.png}

}

\caption{\label{fig-013}Запуск скрипта tangle.jl для генерации форматов
из литературного кода}

\end{figure}%

При выполнении скрипта \texttt{tangle.jl} создались три файла:
\texttt{sir\_literode.ipynb} (Jupyter Notebook),
\texttt{sir\_literode.qmd} (документ Quarto) и \texttt{sir\_literode.jl}
(чистый код). Запускаю файл \texttt{.ipynb} через терминал (рис.
Рисунок~\ref{fig-014}).

\begin{figure}

\centering{

\includegraphics[width=0.7\linewidth,height=\textheight,keepaspectratio]{image/14.png}

}

\caption{\label{fig-014}Запуск сгенерированного Jupyter Notebook}

\end{figure}%

Открываю файл в Jupyter Notebook --- всё работает корректно (рис.
Рисунок~\ref{fig-015}). Можно выполнять ячейки по порядку и видеть
результаты.

\begin{figure}

\centering{

\includegraphics[width=0.7\linewidth,height=\textheight,keepaspectratio]{image/15.png}

}

\caption{\label{fig-015}Содержимое sir\_literode.ipynb в Jupyter
Notebook}

\end{figure}%

Проверяю файл \texttt{sir\_literode.qmd} (рис. Рисунок~\ref{fig-016}).
Он содержит как текст, так и исполняемые ячейки кода, что позволяет в
дальнейшем интегрировать его в итоговый отчёт.

\begin{figure}

\centering{

\includegraphics[width=0.7\linewidth,height=\textheight,keepaspectratio]{image/16.png}

}

\caption{\label{fig-016}Содержание файла sir\_literode.qmd}

\end{figure}%

Также просматриваю чистый код в \texttt{sir\_literode.jl} (рис.
Рисунок~\ref{fig-017}). Он представляет собой извлеченный из
литературного источника исполняемый код без комментариев.

\begin{figure}

\centering{

\includegraphics[width=0.7\linewidth,height=\textheight,keepaspectratio]{image/17.png}

}

\caption{\label{fig-017}Содержание файла с чистым кодом
sir\_literode.jl}

\end{figure}%

\section{Модель
Лотки-Вольтерры}\label{ux43cux43eux434ux435ux43bux44c-ux43bux43eux442ux43aux438-ux432ux43eux43bux44cux442ux435ux440ux440ux44b-1}

После завершения работы с моделью SIR, перехожу к реализации модели
Лотки-Вольтерры.

Создаю файл \texttt{lv\_ode.jl} (рис. Рисунок~\ref{fig-018}).

\begin{figure}

\centering{


\includegraphics[width=0.7\linewidth,height=\textheight,keepaspectratio]{image/18.png}

}

\caption{\label{fig-018}Создание файла lv\_ode.jl для модели
Лотки-Вольтерры}

\end{figure}%

Вставляю программный код, реализующий модель «хищник-жертва» (рис.
Рисунок~\ref{fig-019}). Код содержит функцию \texttt{lotka\_volterra!},
описывающую систему уравнений, а также задаёт классические параметры и
начальные условия.

\begin{figure}

\centering{

\includegraphics[width=0.7\linewidth,height=\textheight,keepaspectratio]{image/19.png}

}

\caption{\label{fig-019}Программный код модели Лотки-Вольтерры в файле
lv\_ode.jl}

\end{figure}%

Запускаю скрипт \texttt{lv\_ode.jl} (рис. Рисунок~\ref{fig-020}).
Система компилирует необходимые пакеты и начинает вычисления.

\begin{figure}

\centering{

\includegraphics[width=0.7\linewidth,height=\textheight,keepaspectratio]{image/20.png}

}

\caption{\label{fig-020}Запуск скрипта lv\_ode.jl}

\end{figure}%

После запуска в консоль выводятся рассчитанные параметры модели,
начальные условия, стационарные точки, а также результаты
статистического и спектрального анализа (рис. Рисунок~\ref{fig-021}).
Рассчитанная стационарная точка ( x\^{}* = 30.0, y\^{}* = 5.0 )
совпадает с теоретическими значениями ( \gamma/\delta ) и (
\alpha/\beta ). Период колебаний жертв составил около 4 единиц времени.

\begin{figure}

\centering{

\includegraphics[width=0.7\linewidth,height=\textheight,keepaspectratio]{image/21.png}

}

\caption{\label{fig-021}Консольный вывод результатов моделирования
Лотки-Вольтерры}

\end{figure}%

На рис. Рисунок~\ref{fig-022} представлен график производных популяций (
dx/dt ) и ( dy/dt ). Он показывает скорость изменения численности
каждого вида во времени.

\begin{figure}

\centering{

\includegraphics[width=0.7\linewidth,height=\textheight,keepaspectratio]{image/22.png}

}

\caption{\label{fig-022}Производные популяций ( dx/dt ) и ( dy/dt )}

\end{figure}%

График «Модель Лотки-Вольтерры: Динамика популяций» (рис.
Рисунок~\ref{fig-023}) является ключевым для анализа. На нём изображены
колебания численности жертв (зелёная линия) и хищников (красная линия).
Хорошо виден классический результат модели --- циклические колебания со
сдвигом фаз: пик численности хищников наступает после пика численности
жертв. Горизонтальными линиями отмечены стационарные уровни ( x\^{}* ) и
( y\^{}* ).

\begin{figure}

\centering{


\includegraphics[width=0.7\linewidth,height=\textheight,keepaspectratio]{image/23.png}

}

\caption{\label{fig-023}Динамика популяций в модели Лотки-Вольтерры}

\end{figure}%

На рис. Рисунок~\ref{fig-024} представлена таблица с результатами,
которая, судя по всему, содержит данные для дальнейшего анализа или была
сгенерирована в процессе отладки.

\begin{figure}

\centering{


\includegraphics[width=0.7\linewidth,height=\textheight,keepaspectratio]{image/24.png}

}

\caption{\label{fig-024}Таблица с результатами моделирования}

\end{figure}%

Фазовый портрет системы (рис. Рисунок~\ref{fig-025}) представляет собой
замкнутую траекторию (предельный цикл) на плоскости ( (x, y) ). Чёрная
точка --- стационарное состояние ( (x\^{}\emph{, y\^{}}) ). Пунктирными
линиями показаны изоклины --- линии, где производная соответствующей
переменной равна нулю. Пересечение изоклин даёт стационарную точку.

\begin{figure}

\centering{

\includegraphics[width=0.7\linewidth,height=\textheight,keepaspectratio]{image/25.png}

}

\caption{\label{fig-025}Фазовый портрет системы «хищник-жертва»}

\end{figure}%

График «Относительные темпы роста» (рис. Рисунок~\ref{fig-026})
показывает скорость изменения популяций в процентах от их текущей
численности ( (dx/dt)/x ) и ( (dy/dt)/y ). Это позволяет сравнивать
динамику видов независимо от масштаба их популяций.

\begin{figure}

\centering{

\includegraphics[width=0.7\linewidth,height=\textheight,keepaspectratio]{image/26.png}

}

\caption{\label{fig-026}Относительные темпы роста популяций}

\end{figure}%

Спектральный анализ (Фурье) колебаний представлен на рис.
Рисунок~\ref{fig-027}. Пики на графике соответствуют основным частотам
колебаний жертв и хищников. Доминирующая частота позволяет определить
основной период колебаний системы.

\begin{figure}

\centering{

\includegraphics[width=0.7\linewidth,height=\textheight,keepaspectratio]{image/27.png}

}

\caption{\label{fig-027}Спектральный анализ (Фурье) колебаний
численности жертв и хищников}

\end{figure}%

После успешного выполнения базового скрипта, я приступил к созданию
литературного кода для модели Лотки-Вольтерры. На рис.
Рисунок~\ref{fig-028}, Рисунок~\ref{fig-029} и Рисунок~\ref{fig-030}
представлены различные версии и фрагменты литературного кода
\texttt{lv\_literate.jl}, где теория сопровождается исполняемыми
блоками.

\begin{figure}

\centering{

\includegraphics[width=0.7\linewidth,height=\textheight,keepaspectratio]{image/28.png}

}

\caption{\label{fig-028}Фрагмент литературного кода модели
Лотки-Вольтерры (версия 1)}

\end{figure}%

\begin{figure}

\centering{

\includegraphics[width=0.7\linewidth,height=\textheight,keepaspectratio]{image/29.png}

}

\caption{\label{fig-029}Фрагмент литературного кода модели
Лотки-Вольтерры (версия 2)}

\end{figure}%

\begin{figure}

\centering{

\includegraphics[width=0.7\linewidth,height=\textheight,keepaspectratio]{image/30.png}

}

\caption{\label{fig-030}Фрагмент литературного кода модели
Лотки-Вольтерры (версия 3)}

\end{figure}%

С помощью скрипта \texttt{tangle.jl} я сгенерировал из литературного
кода чистый код, Jupyter Notebook и Quarto-документацию (рис.
Рисунок~\ref{fig-031}). Все форматы были успешно созданы.

\begin{figure}

\centering{

\includegraphics[width=0.7\linewidth,height=\textheight,keepaspectratio]{image/31.png}

}

\caption{\label{fig-031}Генерация производных форматов для модели
Лотки-Вольтерры}

\end{figure}%

Запускаю сгенерированный Jupyter Notebook \texttt{lv\_literate.ipynb}
(рис. Рисунок~\ref{fig-032}). Сервер Jupyter запускается, и я могу
открыть файл в браузере для интерактивной работы.

\begin{figure}

\centering{

\includegraphics[width=0.7\linewidth,height=\textheight,keepaspectratio]{image/32.png}

}

\caption{\label{fig-032}Запуск Jupyter Notebook для модели
Лотки-Вольтерры}

\end{figure}%

Файл Jupyter Notebook открыт и готов к выполнению (рис.
Рисунок~\ref{fig-033}). Все ячейки с кодом и пояснениями корректно
отображаются.

\begin{figure}

\centering{

\includegraphics[width=0.7\linewidth,height=\textheight,keepaspectratio]{image/33.png}

}

\caption{\label{fig-033}Содержимое файла lv\_literate.ipynb в Jupyter
Notebook}

\end{figure}%

\chapter{Выводы}\label{ux432ux44bux432ux43eux434ux44b}

В ходе выполнения лабораторной работы были успешно реализованы две
классические модели: эпидемиологическая модель SIR и экологическая
модель Лотки-Вольтерры.

Для модели SIR были получены и проанализированы графики динамики
популяций, числа инфицированных, эффективного репродуктивного числа.
Подтверждено, что пик эпидемии наступает при ( R\_e = 1 ), а порог
коллективного иммунитета составляет ( 1-1/R\_0 ). Экспоненциальный рост
на начальном этапе был верифицирован с помощью графика в логарифмическом
масштабе.

Для модели Лотки-Вольтерры получены циклические решения, соответствующие
теоретическим представлениям. Рассчитанные стационарные точки совпали с
ожидаемыми значениями ( x\^{}* = \gamma/\delta ) и ( y\^{}* =
\alpha/\beta ). Фазовый портрет и спектральный анализ подтвердили
колебательный характер системы и позволили определить период колебаний.

Код для обеих моделей был оформлен в литературном стиле с использованием
Literate.jl, что позволило сгенерировать из одного исходного файла
чистый код, Jupyter Notebook и документ Quarto. Это демонстрирует
преимущества подхода литературного программирования для воспроизводимых
научных исследований.

Все поставленные задачи выполнены, полученные результаты визуализированы
и проанализированы.


\printbibliography



\end{document}
